\section*{第 \ref{ch:b2:eruptive-styles} 章——火山喷发}

\begin{solution}{exo:b2:eruptive-styles:1}
(a) $\qty{1.3}{km^3} = \qty{1.3e9}{m^3}$，介于 $10^9$ 与 $10^{10}$ 之间：VEI~5。
(b) $\qty{2.9e8}{m^3}$：VEI~4。(c) $\qty{2e6}{m^3}$，介于 $10^6$ 与 $10^7$ 之间：
VEI~2。(d) VEI~6 起始于 \qty{e10}{m^3}，故\omterm{def:b2:eruptive-styles:tephra}{火山灰碎屑}总量至少须为
\qty{10}{km^3}；“超过 \qty{5}{km^3}”是下界，并非矛盾。
\end{solution}

\begin{solution}{exo:b2:eruptive-styles:2}
(a) 火山灰（细于 \qty{2}{mm}）。(b) \omterm{def:b2:eruptive-styles:tephra}{火山砾}（\qtyrange{2}{64}{mm}）。
(c) 火山块：大于 \qty{64}{mm}，有棱角，喷出时为固态。(d) \omterm{def:b2:eruptive-styles:tephra}{火山弹}，
即纺锤状\omterm{def:b2:eruptive-styles:tephra}{火山弹}：大于 \qty{64}{mm}，
在飞行途中仍为熔融状态时成形。
\end{solution}

\begin{solution}{exo:b2:eruptive-styles:3}
(a) $0.304\times(10^6)^{0.241} = 0.304\times27.9 = \qty{8.5}{km}$；
$0.304\times(10^8)^{0.241} = 0.304\times84.7 = \qty{25.8}{km}$。
(b) $Q = (35/0.304)^{1/0.241} = 115.1^{4.149} = \qty{3.6e8}{kg/s}$，且 $I =
8.55 + 3 = 11.6$。(c) $Q \propto H^{4.15}$，因此高度偏高 \qty{20}{\percent}
会使 $Q$ 被高估 $1.2^{4.15} = 2.1$ 倍。
\end{solution}

\begin{solution}{exo:b2:eruptive-styles:4}
$V_{\mathrm{DRE}} = 0.29\times1000/2500 = \qty{0.116}{km^3}$；$m =
\num{0.29e9}\times1000 = \qty{2.9e11}{kg}$；$M = \log_{10}(\num{2.9e11}) - 7 = 11.46
- 7 = 4.46$。
\end{solution}

\begin{solution}{exo:b2:eruptive-styles:5}
$k = \ln(40/10)/(15 - 5) = 1.386/10 = \qty{0.139}{km^{-1}}$；$T_0 = 40\,e^{0.139\times5}
= 40\times2.00 = \qty{80}{cm}$；$\ln2/k = \qty{5.0}{km}$。取 $T_0 =
\qty{8.0e-4}{km}$：$V = 2\times\num{8.0e-4}/0.139^2 = \qty{0.083}{km^3} =
\qty{8.3e7}{m^3}$，VEI~3。
\end{solution}

\begin{solution}{exo:b2:eruptive-styles:6}
(a) $\int_0^x T_0e^{-ky}2y\,dy = 2T_0\big[-\tfrac{y}{k}e^{-ky} - \tfrac{1}{k^2}e^{-ky}\big]_0^x
= \tfrac{2T_0}{k^2}[1 - (1 + kx)e^{-kx}]$。(b) 迭代 $u = \ln(2(1 + u))$，
取初值 $u = 1$：$1.386$、$1.564$、$1.635$、$1.662$、$1.672$、\dots，收敛于
$u = 1.678$，此时 $(1 + u)e^{-u} = 2.678\times0.1867 = 0.500$。(c) $x_{1/2} =
1.678/0.139 = \qty{12.1}{km}$，位于最外层等厚线（$\sqrt A =
\qty{15}{km}$）之内；然而该等厚线之外的体积占比为 $(1 + 2.09)e^{-2.09} =
\qty{38}{\percent}$。近五分之二的沉积物位于
无任何测量的区域，其体积依赖于外推。
\end{solution}

\begin{solution}{exo:b2:eruptive-styles:7}
(a) $p_s = (5.0/0.444)^2 = \qty{127}{MPa}$，对应深度 $127/0.025 = \qty{5.1}{km}$。
(b) $x = (5.0 - 0.444\sqrt{10})/100 = (5.0 - 1.404)/100 = 0.0360$；
$V_g = 0.0360\times8.314\times1123/(0.018015\times10^7) = \qty{1.86e-3}{m^3}$；
$V_m = 0.964/2300 = \qty{4.19e-4}{m^3}$；$\phi = 1.86/(1.86 + 0.419) = 0.82$。
(c) 玄武质熔体的黏度比流纹质熔体低数千至数十亿倍
（见\cref{ch:b2:magma-emplacement}）：气泡上升、
合并并逸出，气体经由可渗透的泡沫层或火山口排出，
因而\omterm{def:b2:magma-emplacement:magma}{岩浆}以开放系统的方式脱气，而非保留其气体；其
\omterm{def:b2:eruptive-styles:regimes}{岩浆破碎}（若发生）仅限于火山口附近的喷泉区。较慢的上升
为此提供了更多时间。
\end{solution}

\begin{solution}{exo:b2:eruptive-styles:8}
(a) $u_z(0) = 0.75\times10^7/(\pi\times(5000)^2) = \qty{0.095}{m} = \qty{95}{mm}$。
(b) $z_0 = 3.8/0.766 = \qty{5.0}{km}$。(c) 此时为最大值的 $(1 + 1)^{-3/2} = 0.35$
倍。
\end{solution}

\begin{solution}{exo:b2:eruptive-styles:9}
最远射程出现在 \ang{45}：$R = v^2/g = 100^2/9.81 = \qty{1.0}{km}$，
而 $200^2/9.81 = \qty{4.1}{km}$。空气阻力缩短了飞行时间，因此这些结果
是上限，对大块体接近实际，对小块则偏大甚多。
由于 $R = v^2/g$，危险区范围与抛射速度的平方成正比：
喷发速度加倍，其半径增至四倍。
\end{solution}

\begin{solution}{exo:b2:eruptive-styles:10}
(a) $T = T_0e^{-k_1x}$ 适用于 $x < X$，此后 $T = T_Xe^{-k_2(x - X)}$，$T_X =
T_0e^{-k_1X}$（由连续性），利用 $\int_0^X ye^{-ky}dy = \frac{1}{k^2} - e^{-kX}\big(\frac{X}{k} +
\frac{1}{k^2}\big)$ 及 $\int_X^\infty ye^{-ky}dy = e^{-kX}\big(\frac{X}{k} +
\frac{1}{k^2}\big)$：
近端部分给出 $2T_0\big[\frac{1}{k_1^2} - e^{-k_1X}\frac{k_1X + 1}{k_1^2}\big]$，远端
部分为 $2T_Xe^{k_2X}\int_X^\infty ye^{-k_2y}dy = 2T_X\frac{k_2X + 1}{k_2^2}$；代入 $T_X =
T_0e^{-k_1X}$ 后，二者之和即为方法框中的公式。
(b) $T_0 = \qty{0.01285}{km}$：$2T_0/k_1^2 = 0.0257/0.1332 = \qty{0.193}{km^3}$；
$e^{-0.365\times8.56} = 0.0440$；$(0.138\times8.56 + 1)/0.138^2 = 114.5$ 及
$(0.365\times8.56 + 1)/0.365^2 = 30.96$；修正量为 $0.0257\times0.0440\times83.6
= \qty{0.094}{km^3}$，而 $V = \qty{0.287}{km^3}$，与已发表的 \qty{0.29}{km^3} 相符。
$X$ 之外的部分为 $2T_X(X/k_2 + 1/k_2^2) = 2\times\num{5.65e-4}\times114.5 =
\qty{0.129}{km^3}$，即 \qty{45}{\percent}。
\end{solution}

\begin{solution}{exo:b2:eruptive-styles:11}
(a) $100^{1/4} = 3.16$；经验规律给出 $100^{0.241} = 3.03$ 与
$100^{0.226} = 2.83$。(b) $H \propto N^{-3/4}$：层结更强的大气使烟柱
停在更低处；若 $N$ 增大 \qty{20}{\percent}，
则 $1.2^{-0.75} = 0.87$，高度降低 \qty{13}{\percent}。(c) 弱羽流
被风吹弯而降低，同时潮湿空气在其中释放潜热，
使其升高；这些效应对小型与大型喷发的影响不同，
且不包含在简单规律之中。另外，喷发率仅能精确到
数倍范围内，而横轴变量的误差会使
拟合斜率偏向零。
\end{solution}

\begin{solution}{exo:b2:eruptive-styles:12}
(a) 1900 年以前每世纪 $132/400\times100 = 33$ 次；1900 年以后每
世纪 $95/125\times100 = 76$ 次。(b) $33/76 = \qty{43}{\percent}$。(c) $21/56 = \qty{38}{\percent}$：
同样的漏记，相差两至三倍。(d) 较大的喷发
能从更远处被观测到，并留下更厚、更宽、抗侵蚀时间更长的沉积物；
1700 年偏远岛屿上的 VEI~4 喷发可能未被记录，而
VEI~6 喷发则几乎不可能。
\end{solution}

\begin{solution}{pb:b2:eruptive-styles:1}
\textbf{1.} $\ln T = 5.298$，$4.605$，$3.871$，$2.639$，$1.386$。相邻等厚线之间的
斜率分别为 $-0.0693$、$-0.0367$、$-0.0308$ 和 $-0.0313$ 每千米：
一段较陡，位于 10 至 \qty{20}{km} 之间；另一段较缓，位于 40 至
\qty{120}{km}；转折点位于 20 至 \qty{40}{km} 之间。

\textbf{2.} $k_1 = \ln(200/100)/10 = \qty{0.0693}{km^{-1}}$；$T_0 = 200\,e^{0.693} =
\qty{400}{cm}$。

\textbf{3.} $167e^{-1.244} = \qty{48.1}{cm}$，$167e^{-2.488} = \qty{13.9}{cm}$，
$167e^{-3.732} = \qty{4.00}{cm}$：与 48、14 和 4.0 的偏差均在 \qty{1}{\percent} 以内。

\textbf{4.} 由 $400e^{-0.0693X} = 167e^{-0.0311X}$ 可得 $X = \ln(400/167)/(0.0693 -
0.0311) = 0.873/0.0382 = \qty{22.9}{km}$。

\textbf{5.} $\ln2/0.0693 = \qty{10.0}{km}$ 和 $\ln2/0.0311 = \qty{22.3}{km}$。

\textbf{6.} $T_0 = \qty{0.0040}{km}$：$V = 2\times0.0040/0.0693^2 =
\qty{1.67}{km^3}$。

\textbf{7.} $e^{-0.0693\times22.9} = 0.205$；$(0.0311\times22.9 + 1)/0.0311^2 = 1770$；
$(0.0693\times22.9 + 1)/0.0693^2 = 539$；修正项 $2\times0.0040\times0.205\times1231 =
\qty{2.02}{km^3}$；$V = 1.67 + 2.02 = \qty{3.68}{km^3}$，其中修正项占
\qty{55}{\percent}。

\textbf{8.} $3.68/1.67 = 2.2$。远端细粒火山灰层薄，易被侵蚀、
再搬运，或落入海中；在古老的沉积物中，它常常已丢失或从未被圈定，
因此体积被低估两倍或更多。

\textbf{9.} $\qty{3.68e9}{m^3}$ 介于 $10^9$ 与 $10^{10}$ 之间：VEI~5。

\textbf{10.} $V_{\mathrm{DRE}} = 3.68\times900/2400 = \qty{1.38}{km^3}$。

\textbf{11.} $m = \num{3.68e9}\times900 = \qty{3.31e12}{kg}$。

\textbf{12.} $M = \log_{10}(\num{3.31e12}) - 7 = 12.52 - 7 = 5.52$。

\textbf{13.} $3.68/1.3 = 2.8$ 倍。

\textbf{14.} \omterm{def:b2:eruptive-styles:column}{喷发柱}高于 \qty{25}{km}，为\omterm{def:b2:eruptive-styles:styles}{普林尼式喷发}。

\textbf{15.} $Q = \num{3.31e12}/36000 = \qty{9.2e7}{kg/s}$；$I = 7.96 + 3 = 11.0$。

\textbf{16.} $0.304\times(\num{9.2e7})^{0.241} = 0.304\times83.0 = \qty{25.2}{km}$。

\textbf{17.} $0.345\times(\num{9.2e7})^{0.226} = 0.345\times63.1 = \qty{21.8}{km}$。

\textbf{18.} 实测的 \qty{31}{km} 比第一个预测值高 \qty{23}{\percent}，
比第二个预测值高 \qty{42}{\percent}，均在
这些经验规律 \qtyrange{53}{60}{\percent} 的散布范围之内：结果一致。

\textbf{19.} $Q = (31/0.304)^{1/0.241} = 102.0^{4.149} = \qty{2.2e8}{kg/s}$，
为平均速率的 2.3 倍。碎屑记录的是喷发高峰期的最高\omterm{def:b2:eruptive-styles:column}{喷发柱}，
而第~15~题对整个阶段取平均；此外，部分质量
可能未保留在沉积物中（远端火山灰、\omterm{def:b2:eruptive-styles:pdc}{火山碎屑流}）。

\textbf{20.} 是：\omterm{def:b2:eruptive-styles:column}{喷发柱}高达 22 至 \qty{31}{km}，远高于 \qty{16}{km} 处的对流层顶。
进入平流层的气体会扩散到全球；起关键作用的气体是
二氧化硫，皮纳图博火山在 1991 年送入平流层的两千万吨
使全球气温在两年内降低了约半摄氏度。

\textbf{21.} $\lambda_3 = 399/75 = 5.32$，$\lambda_4 = 95/125 = 0.760$，$\lambda_5 = 22/225 =
0.0978$，$\lambda_6 = 8/525 = 0.0152$ 每年；对数分别为 $0.726$、$-0.119$、$-1.010$、
$-1.817$。

\textbf{22.} 拟合直线给出 $10^{0.73} = 5.37$，$10^{-0.12} = 0.759$，$10^{-0.97} =
0.107$ 和 $10^{-1.82} = 0.0151$：与计数相比，偏差分别在 \qty{1}{\percent}、\qty{0}{\percent}、
\qty{10}{\percent} 和 \qty{1}{\percent} 以内。

\textbf{23.} 大喷发的记录在更早的年代仍然完整，
因为它们能从更远处被察觉，沉积物也保存得更久；若
时间段始于 1500 年，则会严重低估 VEI~3 喷发的次数，而
始于 1950 年的时间段则几乎不含 VEI~6 喷发。

\textbf{24.} $\lambda_5 = \num{0.107}$ 每年，重现期为 $1/0.107 =
\num{9.3}\text{ 年}$；十年内发生的概率为 $1 - e^{-1.07} = 1 - 0.343 = \qty{66}{\percent}$。

\textbf{25.} $\lambda_7 = 10^{3.28 - 5.95} = 10^{-2.67} = 0.00214$ 每年，即每
470 年一次；记录给出 $2/1025 = 0.00195$ 每年：两者相符，偏差在
\qty{10}{\percent} 以内，该点未用于拟合。

\textbf{26.} $\lambda_6 = 10^{3.28 - 5.10} = 10^{-1.82} = 0.0151$ 每年；$\lambda t = 1.51$，
$1 - e^{-1.51} = 1 - 0.220 = 0.780$：未来一百年内，全球任一处至少发生一次
VEI~6 或更大喷发的概率为
\textbf{\qty[text-series-to-math=true]{78}{\percent}}。其依据是恒定的速率，
各次喷发相互独立（泊松过程），以及
在各时间段内目录均完整。
\end{solution}
